Misalkan $\gamma$ suatu lintasan-separuh di $M$ seperti yang diberikan, dan misalkan

\mavergleichskette
{\vergleichskette
{ P
}
{ \in }{  U
}
{ \subseteq }{ M
}
{    }{
}
{   }{
}
}
{}{}{}
merupakan suatu daerah peta dengan peta
\maabbdisp {\alpha} { U } {V \subseteq  H_{\leq 0}
} {}
dengan

\mavergleichskettedisp
{\vergleichskette
{  H_{\leq 0}
}
{ =} { \left\{ (x_1,x_2 , \ldots, x_n)  \mid   x_1 \leq 0        \right\}
}
{ } {
}
{ } {
}
{ } {
}
}
{}{}{}
dan

\mavergleichskettedisp
{\vergleichskette
{ Q
}
{ =} { \alpha(P)
}
{ =} { \left(  0   , \,  a_2  , \,   \ldots , \,  a_n                \right)
}
{ } {
}
{ } {
}
}
{}{}{.}
Di bawah pemetaan singgung,

\mavergleichskette
{\vergleichskette
{ v
}
{ = }{ \gamma'(0)
}
{  }{
}
{    }{
}
{   }{
}
}
{}{}{}
dipetakan ke vektor turunan dari

\mavergleichskette
{\vergleichskette
{ \delta
}
{ = }{ \alpha \circ \gamma
}
{  }{
}
{    }{
}
{   }{
}
}
{}{}{}
di titik nol. Karena $\gamma$ bernilai di $M$, $\delta$ bernilai seluruhnya di $H_{\leq 0}$. Oleh karena itu, dalam hasil bagi selisih

\mathdisp {{ \frac{   \left(  \delta_1(t)-\delta_1(0)   , \,   \delta_2(t)-\delta_2(0)  , \,   \ldots , \,   \delta_n(t)-\delta_n(0)                \right)  }{ t } }} {  }

nilai $t$ negatif dan nilai $\delta_1(t)$ nonpositif, sehingga

\mavergleichskettedisp
{\vergleichskette
{  \delta_1'(0)
}
{ \geq} { 0
}
{ } {
}
{ } {
}
{ } {
}
}
{}{}{,}
dan dengan demikian termasuk separuh positif.

Sebaliknya, andaikan
\mathl{T_P(\alpha)(v)}{} termasuk separuh positif. Maka garis
\maabbeledisp {\delta} {\R} {\R^n
} { t } { Q+tT_P(\alpha)(v)
} {}
untuk

\mavergleichskette
{\vergleichskette
{ t
}
{ \leq }{  0
}
{  }{
}
{    }{
}
{   }{
}
}
{}{}{}
terletak seluruhnya di dalam separuh negatif. Lintasan
\mathl{\alpha^{-1} \circ \delta}{} terdefinisi pada suatu interval di sisi negatif, bernilai di $M$, dan merupakan realisasi lintasan-separuh dari $v$.

 [[Kategorie:Latexseite]]
